\FloatBarrier
\section{Discussion}
\label{sec:discussion}

Weighted mixture gives the strongest selected pipeline on held-out data. Random and
weighted use the same architecture, training horizon, mean-$N=5$ inference, split,
seeds, and evaluator; their contrast therefore matches inference compute. Their
training policies differ jointly in source families, size sampling, transforms,
whether current-case sources are eligible, epoch refresh, and fallback. Evidence
supports the full policy, not its configured 80/15/5 proportions alone. Gains in
most paired cases argue against an outlier-only effect, but 15 all-metric losses show
it is not uniformly superior. One training seed and arm-specific checkpoints
prevent component-level causal attribution.

The fixed-mask pipeline uses $N=1$ by design. On development data, mean $N=3$
improved joint rank by only 0.1333 at three times the trajectory count, while mean
$N=5$ was worse. Thus fixed versus augmented compares unequal trajectory counts;
weighted versus random is the equal-trajectory contrast.

The invertible wavelet representation enables full-volume 3D sampling, while hard
compositing preserves observed tissue. Sampling remains costly: mean $N=5$ runs
five 1000-step ancestral trajectories. Recent work on faster wavelet
inpainting~\cite{durrer2026fastwdm3d} motivates sampler validation under the same
frozen cohort and scoring protocol.

Metrics cover artificially removed healthy voxels. Because observed T1n remains
pathological in $m_u$, neither the metrics nor the qualitative reference establish
patient-specific counterfactual correctness there; such ground truth does not
exist. Plausible reconstructions may hallucinate incorrect anatomy and are not
clinically validated. Hidden organizer-validation ground truth prevents local
external scoring.

The case-keyed split is not fully patient-independent: optimization includes sibling
timepoints for 3/25 development and 11/75 confirmation cases; development and
confirmation patients are mutually disjoint. In a post-hoc 64-case patient-disjoint
analysis, weighted mixture obtains SSIM 0.8101, PSNR 19.27\,dB, and MSE 0.009220;
random obtains 0.7911, 18.33\,dB, and 0.011664; and fixed obtains 0.7703,
17.93\,dB, and 0.012764. The ordering is unchanged. This removes confirmation
overlap only, not possible influence on development-stage selection.

Other limitations are one seed per variant, five-case one-trajectory checkpoint
diagnostics, transfer of terminally selected inference policies to arm-specific
checkpoints, and no completed external baseline. Across-case variability does not
measure seed variability, so prior-method superiority is unestablished. Development
ranks lack their original per-case predictions; one random-arm dirty-worktree diff
was not preserved, although its checkpoint, configuration, metrics, sampling
metadata, and prediction hashes are archived. Future work should add seeds and
baselines, preserve all outputs, validate faster sampling, and use patient-level
splits and multimodal conditioning.

\section{Conclusion}
\label{sec:conclusion}

CATCH combines conditional 3D Haar-wavelet diffusion, tumor-excluded supervision,
hard preservation of observed tissue, and development-selected trajectory
aggregation. In this single-seed internal study, weighted-mixture augmentation is
strongest under matched compute on artificially removed healthy tissue; clinical
and cross-method superiority remain unestablished.
